\documentclass[11pt]{article}

% ============================================================
% EV-SUM-03 — Evaluación de la unidad 3
% Ingeniería de Requerimientos (ISR-401) · UTEQ · SPA 2026-2027
% Compilación: pdflatex -> bibtex -> pdflatex -> pdflatex
% ============================================================

\usepackage[utf8]{inputenc}
\usepackage[T1]{fontenc}
% babel no disponible en este entorno
\usepackage[scaled]{helvet}
\renewcommand{\familydefault}{\sfdefault}
\usepackage[a4paper,top=2.6cm,bottom=2.3cm,left=2cm,right=2cm,headheight=34pt]{geometry}
\usepackage{graphicx}
\usepackage{array}
\usepackage{tabularx}
\usepackage{longtable}
\usepackage{pdflscape}
\usepackage{colortbl}
\usepackage{xcolor}
\usepackage{enumitem}
\usepackage{fancyhdr}
\usepackage{lastpage}
\usepackage[most]{tcolorbox}
\usepackage{amsmath}
\usepackage{amssymb}
\usepackage[numbers,sort&compress]{natbib}
\usepackage{url}
\usepackage{needspace}
\usepackage{caption}
\usepackage[hidelinks]{hyperref}

% ---------- Textos en español (sin babel) ----------
\renewcommand{\tablename}{Tabla}
\renewcommand{\refname}{Referencias}

% ---------- Paleta institucional ----------
\definecolor{utegreen}{HTML}{0B7A3B}
\definecolor{utegreendk}{HTML}{075C2C}
\definecolor{midblue}{HTML}{1F5C8B}
\definecolor{gold}{HTML}{D4A017}
\definecolor{rowgray}{HTML}{EFEFEF}
\definecolor{boxgray}{HTML}{F4F6F4}
\definecolor{lightred}{HTML}{FBEAEA}
\definecolor{lightgold}{HTML}{FBF3DC}
\definecolor{lightblue}{HTML}{EAF1F7}
\definecolor{lightgreen}{HTML}{DCEFE3}

% ---------- Columnas personalizadas ----------
\newcolumntype{L}[1]{>{\raggedright\arraybackslash}p{#1}}
\newcolumntype{C}[1]{>{\centering\arraybackslash}p{#1}}
\newcolumntype{Y}{>{\raggedright\arraybackslash}X}

% ---------- Encabezado y pie ----------
\pagestyle{fancy}
\fancyhf{}
\lhead{\footnotesize UTEQ \,$\cdot$\, Ingeniería de Requerimientos (ISR-401)}
\rhead{\footnotesize EV-SUM-03 \,$\cdot$\, SPA 2026--2027}
\lfoot{\footnotesize Prof. Gleiston C. Guerrero-Ulloa, PhD}
\rfoot{\footnotesize Página \thepage\ de \pageref{LastPage}}
\renewcommand{\headrulewidth}{0.6pt}
\renewcommand{\footrulewidth}{0.4pt}

\setlength{\parindent}{0pt}
\setlength{\parskip}{5pt}
\setlength{\emergencystretch}{3em}
\renewcommand{\arraystretch}{1.25}

% ---------- Atajos ----------
\newcommand{\thdr}[1]{\cellcolor{utegreen}\textcolor{white}{\textbf{#1}}}
\newcommand{\seccion}[1]{\section{\textcolor{utegreendk}{#1}}}
\newcommand{\subseccion}[1]{\subsection{\textcolor{midblue}{#1}}}
\renewcommand{\thesection}{\arabic{section}.}
\renewcommand{\thesubsection}{\arabic{section}.\arabic{subsection}.}

\newtcolorbox{cajaclave}[1][]{colback=boxgray,colframe=utegreen,boxrule=0.8pt,arc=2mm,fonttitle=\bfseries,breakable,#1}
\newtcolorbox{cajapiso}[1][]{colback=lightred,colframe=red!60!black,boxrule=1pt,arc=2mm,fonttitle=\bfseries,breakable,#1}
\newtcolorbox{cajaaviso}[1][]{colback=lightgold,colframe=gold!80!black,boxrule=0.8pt,arc=2mm,fonttitle=\bfseries,breakable,#1}

\begin{document}

\noindent\includegraphics[width=\linewidth]{logoband.png}

\vspace{4pt}
\begin{center}
	{\small Facultad de Ciencias de la Computación \,\textcolor{gold}{$\bullet$}\, Carrera de Software \,\textcolor{gold}{$\bullet$}\, Ingeniería de Requerimientos (ISR-401)}\\[8pt]
	{\LARGE\textbf{\textcolor{utegreendk}{GUÍA Y RÚBRICA DE EVALUACIÓN}}}\\[4pt]
	{\large\textbf{\textcolor{utegreendk}{ACTIVIDAD SUMATIVA}}}\\[6pt]
	{\Large\textbf{EV-SUM-03}}\\[2pt]
	{\large Evaluación de la unidad 3}
\end{center}

\vspace{2pt}
\noindent\textcolor{gold}{\rule{\linewidth}{1.5pt}}

\seccion{Datos generales}

\noindent\begin{tabularx}{\linewidth}{|L{0.30\linewidth}|Y|}
	\hline
	\thdr{Campo} & \thdr{Detalle}\\
	\hline
	Asignatura & Ingeniería de Requerimientos (ISR-401) \,$\cdot$\, 4.\textsuperscript{o} nivel \,$\cdot$\, Paralelo A\\
	\hline
	\rowcolor{rowgray} Docente & Gleiston Cicerón Guerrero Ulloa, PhD \,$\cdot$\, \texttt{gguerrero@uteq.edu.ec}\\
	\hline
	Unidad & Unidad 3. Especificación de Requisitos de Software\\
	\hline
	\rowcolor{rowgray} Resultado de aprendizaje & Elabora especificaciones de requisitos (ERS/SRS) completas y verificables aplicando IEEE 29148 y UML.\\
	\hline
	Tipo y componente & Evaluación de unidad\\
	\hline
	\rowcolor{rowgray} Peso en el corte & 1.00 de 10 puntos del corte 2\\
	\hline
	Formativa vinculada (requisito) & FORM-TL-05\\
	\hline
	\rowcolor{rowgray} Modalidad & Individual\\
	\hline
	Semana & 10 (28 de diciembre -- 3 de enero de 2027)\\
	\hline
	\rowcolor{rowgray} Publicación de la guía & lunes 21 de diciembre de 2026\\
	\hline
	Cierre & En clase, al término de la sesión del día de clase de la semana 10 (28 de diciembre -- 3 de enero de 2027), según el horario del paralelo\\
	\hline
	\rowcolor{rowgray} Repositorio & Repositorio personal \texttt{ISR401-<ApellidoNombre>}\\
	\hline
	Escala & 0 a 10 puntos, con un decimal\\
	\hline
\end{tabularx}

\seccion{Propósito}

Comprobar de forma individual el logro del resultado de aprendizaje de la unidad 3, combinando un cuestionario en el SGA (conocimiento) y una prueba práctica sobre un caso nuevo (aplicación). Ambas partes se rinden el mismo día, en la misma sesión.

\seccion{Temario evaluado}

\begin{itemize}[leftmargin=20pt,itemsep=2pt]
	\item Documentación: objetivos, tipos de artefactos, formatos de representación, atributos y criterios de calidad \citep{pohl_2025}.
	\item Documento ERS: tipos de especificaciones, estructura de ISO/\allowbreak{}IEC/\allowbreak{}IEEE 29148:2018, lenguaje natural controlado, plantillas y restricciones sintácticas \citep{iso29148_2018,mavin_ears_2009}.
	\item Combinación de modelos UML y texto; modelos de especificación UML \citep{omg_uml_2017}.
\end{itemize}

\seccion{Estructura de la evaluación}

\noindent\begin{tabularx}{\linewidth}{|L{0.20\linewidth}|Y|C{0.12\linewidth}|C{0.10\linewidth}|}
	\hline
	\thdr{Parte} & \thdr{Descripción} & \thdr{Tiempo} & \thdr{Peso}\\
	\hline
	Q: Cuestionario (SGA) & 20 preguntas de opción múltiple, verdadero o falso con justificación y relacionar. Se registra como formativo en la planificación del SGA; su nota entra en la fórmula de la evaluación. & 25 min & 40\,\%\\
	\hline
	\rowcolor{rowgray} P: Prueba práctica & Caso nuevo entregado al inicio de la sesión; se resuelve en \LaTeX{} en la carpeta \texttt{EV-SUM-03/} del repositorio personal, con commits durante la sesión. & 60 min & 60\,\%\\
	\hline
\end{tabularx}

\seccion{Tareas tipo de la parte práctica}

El enunciado real se entrega el día de la evaluación. Las tareas serán de este tipo:

\begin{enumerate}[leftmargin=20pt,itemsep=2pt]
	\item Identificar y corregir los defectos de 6 requisitos ambiguos usando los criterios de lenguaje de 29148 y plantillas EARS.
	\item Completar los atributos de 4 requisitos dados (prioridad, fuente, estabilidad, criterio de aceptación).
	\item Ubicar 8 fragmentos de información en la sección correcta de la estructura de la ERS.
	\item Elaborar un diagrama de secuencia de un escenario dado y vincularlo con los requisitos que representa.
\end{enumerate}

\seccion{Desarrollo de la actividad}

\begin{enumerate}[leftmargin=20pt,itemsep=4pt]
	\item \textbf{Antes de la sesión.} Verifique que su repositorio personal compila y que puede hacer commits desde el equipo que usará. Tenga lista la plantilla \LaTeX{} de la asignatura.
	\item \textbf{Cuestionario.} Rinda el cuestionario en el SGA en el tiempo asignado. Es individual y sin consulta.
	\item \textbf{Prueba práctica.} Cree la carpeta \texttt{EV-SUM-03/}, resuelva cada tarea en su \texttt{main.tex} y haga al menos un commit cada 20 minutos. Solo se permite consultar la norma ISO/\allowbreak{}IEC/\allowbreak{}IEEE 29148 y los apuntes propios; no se permiten asistentes de IA ni comunicación con otras personas.
	\item \textbf{Cierre.} Compile, suba el PDF y el fuente antes de la hora de cierre y suba al SGA el PDF de una sola página con la URL del repositorio.
	\item \textbf{Revisión.} La evaluación se analiza con los estudiantes en la siguiente sesión, antes de registrar la nota.
\end{enumerate}

\seccion{Productos entregables}

{\small
\begin{longtable}{|L{0.27\linewidth}|L{0.43\linewidth}|L{0.21\linewidth}|}
	\caption{Productos que deben estar en el repositorio al cierre.}\\
	\hline
	\thdr{Ruta en el repositorio} & \thdr{Contenido} & \thdr{Mínimo exigido}\\
	\hline
	\endfirsthead
	\hline
	\thdr{Ruta en el repositorio} & \thdr{Contenido} & \thdr{Mínimo exigido}\\
	\hline
	\endhead
	\texttt{EV-SUM-03/\allowbreak{}main.tex} & Resolución de la parte práctica & Todas las tareas\\
	\hline
	\rowcolor{rowgray} \texttt{EV-SUM-03/\allowbreak{}main.pdf} & PDF compilado & 1\\
	\hline
\end{longtable}
}

\seccion{Normas de entrega}

\begin{enumerate}[leftmargin=20pt,itemsep=3pt]
	\item \textbf{Repositorio.} Todo el trabajo se desarrolla en el repositorio personal \texttt{ISR401-<ApellidoNombre>}, que debe ser público en GitHub. Allí deben estar los fuentes (\texttt{.tex}, \texttt{.bib}, figuras, fuentes editables de los modelos y evidencias) y el PDF compilado.
	\item \textbf{Documento en el SGA.} En la tarea del SGA se sube solamente un PDF de \textbf{una sola página} con los datos informativos: estudiante o estudiantes, código y nombre de la tarea, fecha límite de entrega y la URL del repositorio escrita completa en una sola línea. No se sube ningún otro archivo.
	\item \textbf{\LaTeX{} reproducible.} El PDF debe poder regenerarse clonando el repositorio y compilando el archivo \texttt{.tex}. El \texttt{README.md} documenta compilador, archivo principal, orden de comandos y dependencias.
	\item \textbf{Referencias.} Estilo IEEE. Toda referencia debe existir y sus datos (autores, año, título, revista o editorial, volumen, número, páginas y DOI o ISBN) deben corresponder exactamente a la obra real.
	\item \textbf{Plazo.} La parte práctica se resuelve y se sube al repositorio durante la sesión de clase del día de clase de la semana 10 (28 de diciembre -- 3 de enero de 2027), según el horario del paralelo; el cierre es la hora de finalización de la sesión. El cuestionario se rinde en el SGA en la misma sesión.
	\item \textbf{Datos personales.} Ningún dato personal de terceros (cédulas, teléfonos, correos personales, audios o videos de participantes, consentimientos firmados) se publica en el repositorio. Esas evidencias se guardan en una carpeta institucional privada (cuenta \texttt{@uteq.edu.ec}) compartida con el docente, y en el repositorio se registra un manifiesto con nombre de archivo, fecha, duración y huella SHA-256 \citep{lopdp_2021}.
	\item \textbf{Uso de IA como apoyo.} Si se usa una herramienta de IA solo como apoyo (ortografía, error de compilación), se declara en una sección del documento: herramienta, versión, propósito y cómo se verificó lo obtenido.
\end{enumerate}

\seccion{Criterios de piso}

\begin{cajapiso}[title={Criterios de piso: su incumplimiento implica calificación CERO en la actividad}]
	\begin{enumerate}[label=\textbf{P\arabic*.},leftmargin=24pt,itemsep=3pt]
		\item \textbf{PDF de una sola página con datos informativos y URL.} En el SGA se sube solamente un PDF de UNA SOLA página con los datos informativos (estudiante o estudiantes, tarea, fecha límite de entrega) y, en una sola línea, la URL del repositorio. Si el PDF tiene más de una página, si se suben otros archivos, si la URL no cabe en una sola línea o está rota, o si el repositorio no es accesible, la calificación de la actividad es \textbf{CERO}.
		\item \textbf{PDF reproducible.} Si el PDF no se genera correctamente a partir del \LaTeX{} al clonar el repositorio y compilar el archivo \texttt{.tex} siguiendo el \texttt{README.md} (compilador, orden de comandos, archivo principal y dependencias), la calificación de la actividad es \textbf{CERO}.
		\item \textbf{Plazo.} Si el contenido evaluable se sube después de la fecha y hora de cierre, la calificación de la actividad es \textbf{CERO}.
		\item \textbf{Integridad.} Plagio verificado, cualquiera de los indicios graves G1 a G6 de uso de IA, o contenido generado por IA presentado como propio: la calificación de la actividad es \textbf{CERO}, sin perjuicio de las acciones disciplinarias que correspondan.
		\item \textbf{Formativa vinculada.} Quien no entregó de forma válida la actividad formativa vinculada (FORM-TL-05) no puede presentar esta actividad y obtiene CERO en ella, salvo justificación aprobada.
	\end{enumerate}
	En esta actividad, P1 y P2 se aplican al material entregado en el repositorio; la inasistencia no justificada a la sesión implica CERO.
\end{cajapiso}

\seccion{Integridad académica e indicios de uso de IA}

Todo trabajo debe ser de autoría propia. Se aplica el Código de Ética de la ACM \citep{acm_code_2018} y la sección 6.1 del encuadre.

\begin{cajapiso}[title={Indicios graves de uso de IA: cada uno, por sí solo, implica CERO (P4)}]
	\begin{enumerate}[label=\textbf{G\arabic*.},leftmargin=24pt,itemsep=3pt]
		\item \textbf{Referencias con datos que no corresponden a la obra real:} DOI que lleva a una obra con autores, título, año, revista, volumen o páginas diferentes; DOI, ISBN o URL inexistentes; obras que no existen.
		\item \textbf{Citas que la fuente no respalda:} afirmaciones, cifras o citas textuales atribuidas a una fuente que no las contiene, o cláusulas, secciones o tablas de una norma (IEEE 29148, ISO/\allowbreak{}IEC 25010) que no existen o no dicen lo afirmado.
		\item \textbf{Restos de un asistente de IA} en el documento, el código, los comentarios o los commits (por ejemplo, «Claro, aquí tienes\ldots», «Como modelo de lenguaje\ldots»).
		\item \textbf{Evidencias de campo fabricadas:} entrevistas, encuestas, participantes, transcripciones, actas o fotografías que no existen o no tienen respaldo original.
		\item \textbf{Resultados o artefactos inexistentes:} el texto describe modelos, matrices, pruebas o métricas que no están en el repositorio o no coinciden con él.
		\item \textbf{Incapacidad de explicar lo entregado:} el estudiante no puede justificar ni modificar, cuando se le pide, lo que escribió, modeló o especificó.
	\end{enumerate}
\end{cajapiso}

\begin{cajaaviso}[title={Indicios complementarios (se acumulan)}]
	\begin{enumerate}[label=\textbf{I\arabic*.},leftmargin=24pt,itemsep=3pt]
		\item Redacción uniforme de principio a fin, o estilo distinto al del estudiante en otras entregas o en su exposición oral.
		\item Uso excesivo de los mismos conectores y fórmulas («además», «cabe destacar», «es importante señalar», «en conclusión»).
		\item Contenido genérico: requisitos, actores o reglas aplicables a cualquier sistema, sin relación con lo elicitado.
		\item Entrega en uno o pocos commits masivos, o mensajes de commit idénticos o automáticos.
		\item Residuos de formato ajeno a \LaTeX{} (marcas de Markdown, emojis) o metadatos del PDF de otra herramienta.
		\item Estructura, redacción o errores muy similares entre estudiantes o equipos distintos.
		\item Traducción literal o mezcla de idiomas sin justificación.
	\end{enumerate}
	Con dos indicios complementarios se resta 1,5 puntos (nivel B) y con tres o más se restan 3,0 (nivel C). La declaración de uso de IA no exime de responsabilidad ni anula estos indicios.
\end{cajaaviso}

\seccion{Escala de desempeño}

\begin{cajaclave}[title={La presencia no basta: se califica la calidad}]
	Que un indicador esté presente no basta para puntuarlo. Cada indicador se valora con los aspectos de calidad del encuadre: Completitud, veracidad, precisión técnica, claridad y no ambigüedad, verificabilidad, profundidad, pertinencia y relevancia, suficiencia de la evidencia, coherencia y consistencia, trazabilidad, originalidad y autoría, constancia, reproducibilidad, y actualidad y calidad de las fuentes. El nivel asignado a una dimensión es el que corresponde al aspecto peor cumplido entre los que aplican. Un indicador ausente, no verificable en el repositorio o cuyo contenido sea falso vale 0.
\end{cajaclave}

Cada dimensión se califica con un nivel de 0 a 4: 4 Excelente (cubre el 100\,\% con calidad plena), 3 Bueno (al menos 75\,\%, deficiencias menores), 2 Aceptable (50--74\,\% o descriptivo sin análisis), 1 Insuficiente (menos del 50\,\%, superficial o con errores graves) y 0 Ausente. Los descriptores de nivel son los de la rúbrica general del encuadre; la columna de indicadores concreta qué se observa en esta actividad \citep{brookhart_2018,panadero_jonsson_2020}.

\begin{landscape}
\seccion{Rúbrica analítica}
{\footnotesize\renewcommand{\arraystretch}{1.12}
\begin{longtable}{|L{2.4cm}|L{5.6cm}|L{3.1cm}|L{2.8cm}|L{2.8cm}|L{2.7cm}|L{2.1cm}|}
	\caption{Rúbrica analítica de EV-SUM-03. Pesos del tipo de actividad EV según el encuadre.}\\
	\hline
	\thdr{Dimensión (peso)} & \thdr{Indicadores de esta actividad} & \thdr{4 Excelente} & \thdr{3 Bueno} & \thdr{2 Aceptable} & \thdr{1 Insuficiente} & \thdr{0 Ausente}\\
	\hline
	\endfirsthead
	\hline
	\thdr{Dimensión (peso)} & \thdr{Indicadores de esta actividad} & \thdr{4 Excelente} & \thdr{3 Bueno} & \thdr{2 Aceptable} & \thdr{1 Insuficiente} & \thdr{0 Ausente}\\
	\hline
	\endhead
	\hline
	\multicolumn{7}{r}{\footnotesize\textit{Continúa en la página siguiente\ldots}}\\
	\endfoot
	\hline
	\endlastfoot
	\textbf{D1 Cumplimiento y completitud} (10\,\%) & \begin{itemize}[leftmargin=9pt,itemsep=0pt,topsep=0pt,parsep=0pt,partopsep=0pt,after=\vspace{-\baselineskip}]\item Resuelve todas las tareas de la parte práctica en el tiempo asignado.\item Entrega el \texttt{main.tex} y el PDF en la carpeta indicada antes del cierre de la sesión.\end{itemize} & Entrega todos los productos pedidos en la guía, completos, con la estructura exigida y en el repositorio. & Falta un elemento menor o uno está incompleto sin afectar el conjunto. & Faltan varios elementos o hay productos incompletos que limitan la evaluación. & Entrega una parte mínima de lo pedido. & No entrega o lo entregado no corresponde a la actividad.\\
	\hline
	\textbf{D2 Fundamentación conceptual} (30\,\%) & \begin{itemize}[leftmargin=9pt,itemsep=0pt,topsep=0pt,parsep=0pt,partopsep=0pt,after=\vspace{-\baselineskip}]\item Usa con precisión los conceptos y la terminología de la unidad.\item Justifica cada respuesta con el concepto o la norma que la sustenta.\end{itemize} & Define y aplica con precisión los conceptos de la unidad, los relaciona entre sí y los sustenta en la bibliografía y las normas. & Conceptos correctos con alguna imprecisión menor o relación incompleta entre ellos. & Conceptos repetidos de la fuente sin aplicarlos, o con errores que no invalidan el trabajo. & Errores conceptuales graves o confusión entre conceptos clave. & No hay fundamentación o es ajena al tema.\\
	\hline
	\textbf{D3 Desarrollo técnico y evidencia} (30\,\%) & \begin{itemize}[leftmargin=9pt,itemsep=0pt,topsep=0pt,parsep=0pt,partopsep=0pt,after=\vspace{-\baselineskip}]\item Los artefactos pedidos (requisitos, tablas, modelos) son correctos y se aplican al caso del enunciado.\item Los cálculos y clasificaciones son correctos.\end{itemize} & Artefactos completos y correctos; cada afirmación está respaldada por evidencia original verificable en el repositorio. & Artefactos correctos con fallos menores; casi todas las afirmaciones tienen evidencia. & Artefactos parciales o evidencia incompleta para comprobar lo afirmado. & Artefactos incorrectos o la evidencia no permite verificar casi nada. & No hay artefactos ni evidencia verificable.\\
	\hline
	\textbf{D4 Análisis crítico y calidad de los requisitos} (20\,\%) & \begin{itemize}[leftmargin=9pt,itemsep=0pt,topsep=0pt,parsep=0pt,partopsep=0pt,after=\vspace{-\baselineskip}]\item Analiza el caso, decide y justifica; no se limita a definir.\item Los requisitos formulados son claros, singulares y verificables.\end{itemize} & Requisitos claros, verificables, consistentes y trazables; interpreta, compara alternativas, reconoce limitaciones y concluye con afirmaciones sustentadas. & Análisis correcto con alguna conclusión o requisito sin sustento suficiente. & Describe sin interpretar, o hay ambigüedades que afectan la verificación. & Conclusiones o requisitos genéricos que no se derivan de lo realizado. & No hay análisis ni conclusiones.\\
	\hline
	\textbf{D5 Comunicación (documento o exposición)} (10\,\%) & \begin{itemize}[leftmargin=9pt,itemsep=0pt,topsep=0pt,parsep=0pt,partopsep=0pt,after=\vspace{-\baselineskip}]\item Respuestas ordenadas, numeradas por tarea y legibles en el PDF.\item Redacción técnica precisa.\end{itemize} & Documento \LaTeX{} limpio (sin desbordes ni referencias sin resolver), redacción técnica clara; en exposición, dominio del tema, tiempo respetado y respuestas precisas. & Defectos menores de formato, redacción o fluidez. & Varios defectos que dificultan la lectura o la comprensión de la exposición. & Documento desordenado o exposición leída sin dominio del tema. & No hay documento legible o no expone.\\
	\hline
	\textbf{D6 Referencias, ética e integridad} (0\,\%) & \multicolumn{6}{L{20.8cm}|}{No se pondera en las evaluaciones de unidad: la integridad se controla con el criterio de piso P4.}\\
	\hline
\end{longtable}
}
{\footnotesize Nota sobre D6: una referencia inexistente, o cuyo DOI lleve a una obra con datos diferentes a los citados, no se califica en D6: implica calificación CERO (criterio de piso P4).}
\end{landscape}

\seccion{Hoja de cálculo de la nota}

\[
	\text{Nota previa} = \frac{\sum_{i=1}^{6} w_i \times n_i}{4} \times \frac{10}{100}, \qquad \text{Nota final} = \max\left(0,\ \text{Nota previa} - \text{Penalizaciones}\right)
\]

donde $w_i$ es el peso de la dimensión $i$ (en porcentaje) y $n_i$ el nivel obtenido (0 a 4). Si se incumple un criterio de piso, la nota final es cero. La retroalimentación informa siempre la nota previa y la nota final.

\noindent\begin{tabularx}{\linewidth}{|C{0.07\linewidth}|Y|C{0.10\linewidth}|C{0.14\linewidth}|C{0.20\linewidth}|}
	\hline
	\thdr{Dim.} & \thdr{Descripción} & \thdr{Peso (\%)} & \thdr{Nivel (0--4)} & \thdr{$w_i \times n_i$}\\
	\hline
	D1 & Cumplimiento y completitud & 10 &  & \\
	\hline
	\rowcolor{rowgray} D2 & Fundamentación conceptual & 30 &  & \\
	\hline
	D3 & Desarrollo técnico y evidencia & 30 &  & \\
	\hline
	\rowcolor{rowgray} D4 & Análisis crítico y calidad de los requisitos & 20 &  & \\
	\hline
	D5 & Comunicación (documento o exposición) & 10 &  & \\
	\hline
	\rowcolor{rowgray} D6 & Referencias, ética e integridad & 0 & --- & \\
	\hline
	\multicolumn{2}{|r|}{\cellcolor{lightgold}\textbf{Suma}} & \cellcolor{lightgold}\textbf{100} & \cellcolor{lightgold} & \cellcolor{lightgold}\rule{0pt}{12pt}\\
	\hline
	\multicolumn{4}{|r|}{\cellcolor{lightblue}\textbf{Nota previa = Suma / 4 $\times$ 0,1}} & \cellcolor{lightblue}\rule{0pt}{12pt}\\
	\hline
	\multicolumn{4}{|r|}{\cellcolor{lightblue}\textbf{Penalizaciones (tabla siguiente)}} & \cellcolor{lightblue}\rule{0pt}{12pt}\\
	\hline
	\multicolumn{4}{|r|}{\cellcolor{lightgreen}\textbf{Nota final de la actividad (sobre 10)}} & \cellcolor{lightgreen}\rule{0pt}{12pt}\\
	\hline
\end{tabularx}

\begin{cajaclave}[title={Nota de la evaluación de unidad}]
	La tabla anterior califica la parte práctica (P). La nota de la evaluación es $\text{EV} = 0{,}40\,Q + 0{,}60\,P$, donde $Q$ es la nota del cuestionario del SGA sobre 10.
\end{cajaclave}

\subseccion{Penalizaciones que no anulan la nota}

\noindent\begin{tabularx}{\linewidth}{|Y|C{0.18\linewidth}|C{0.12\linewidth}|}
	\hline
	\thdr{Situación detectada} & \thdr{Efecto (sobre 10)} & \thdr{Aplica}\\
	\hline
	Falta la declaración de uso de inteligencia artificial & $-0{,}3$ & \\
	\hline
	\rowcolor{rowgray} Referencias en formato APA u otro distinto de IEEE & D6 = 0 & \\
	\hline
	Dos indicios complementarios de uso de IA (nivel B) & $-1{,}5$ & \\
	\hline
	\rowcolor{rowgray} Tres o más indicios complementarios (nivel C) & $-3{,}0$ & \\
	\hline
\end{tabularx}

\bibliographystyle{unsrtnat}
\bibliography{referencias}

\end{document}
